\chapteropen{05}

\begin{cele}
\item Policzyć okno snu w godzinach i zobaczyć, ile z niego zostaje na sen, a ile na czekanie, budzenie się i przyzwyczajenie.
\item Ustalić stałą porę snu i wstawania, także w weekend, i uzasadnić, dlaczego to właśnie pobudka wyznacza rytm.
\item Rozpisać światło, kofeinę, alkohol i sypialnię na konkretne godziny, a nie na ogólne zasady.
\item Rozpoznać moment, w którym długi czas w łóżku przestaje być zasobem, i przesunąć porę snu o kwadranse zamiast o skoki.
\end{cele}

\section{Stała pora, także w sobotę}

Rytm snu to nie kwestia tego, ile godzin śpisz, tylko tego, o której. Organizm przygotowuje się do snu i do pobudki z wyprzedzeniem, mniej więcej tak jak rozkład jazdy: tory są te same, a skład przyjeżdża o tej samej porze, więc rzadziej coś się zacina. Stała pora jest najtańszym narzędziem w całej tej książce. Nie kosztuje pieniędzy ani sprzętu, kosztuje jedną decyzję i jej konsekwencję.

Najmocniej działa godzina wstawania. Budzik o tej samej porze każdego dnia daje ciału powtarzalny sygnał, a godzina zaśnięcia przesuwa się wtedy sama, w granicach tego, co zmęczenie na to pozwala. \verify{Dla większości dorosłych stała pora wstawania silniej porządkuje rytm dobowy niż stała pora zasypiania.} Jeśli wstajesz o 6:00 w poniedziałek, a o 9:00 w sobotę, to w praktyce dostajesz dwie różne doby, które tylko udają jedną.

Weekend jest miejscem, w którym ta zasada najczęściej pada. Różnica między dniami roboczymi a wolnymi, nazywana czasem społecznym jet lagiem, polega na tym, że w tygodniu śpisz według zegara pracy, a w sobotę według własnej woli. Przykład jest prosty: w tygodniu okno snu to 21:00--09:00, w sobotę 23:00--11:00. Środek snu przesuwa się z trzeciej w nocy na piątą rano, czyli o dwie godziny. Jedna taka sobota to niewiele. Pięćdziesiąt weekendów w roku to już rytm ustawiony na inną strefę niż ta, w której pracujesz. \verify{Skutki zdrowotne regularnego przesuwania środka snu o kilka godzin są przedmiotem badań i nie należy ich traktować jako ustalonych.}

Szablon weekendowy w tej książce zachowuje dwanaście godzin snu, czyli \val{week.C.sleep} godzin, tak jak dni robocze. Zmienia się to, co dzieje się w ciągu doby poza snem, a nie długość samego snu. Przy wolnych dniach pokusa, żeby \glqq{}odespać\grqq{} tydzień, jest największa. Dlatego zasada brzmi: w weekend zmieniasz tyle, ile potrzebujesz, żeby wstawać później o godzinę, a nie o trzy.

\section{Okno 12 godzin, rozpisane na godziny}

Okno to czas między położeniem się do łóżka a wstaniem, a nie czas snu. Różnica jest ważna, bo sen zajmuje tylko część okna. Zasypianie trwa, budzenie się w nocy też, a czasem leżysz i czekasz, zanim zrozumiesz, że nie zaśniesz. Skoro okno wynosi dwanaście godzin, a sen ma trwać około ośmiu (rozdział o potrzebach snu mówi, dlaczego tyle), to reszta okna musi mieć swoje miejsce w budżecie.

\begin{przyklad}[Doba 21:00--09:00, godzina po godzinie]
\begin{itemize}
\item \textbf{20:15--21:00} rutyna wieczorna, 45 minut (szczegóły w punkcie o rutynie).
\item \textbf{21:00} w łóżku, światło przygaszone, telefon poza sypialnią.
\item \textbf{21:30} zakładamy zasypianie po około 30 minutach. \verify{Typowy czas zasypiania u dorosłych mieści się zwykle w przedziale kilkunastu do kilkudziesięciu minut.}
\item \textbf{około 05:00} pół godziny ocknięcia się w nocy, bez wstawania z łóżka.
\item \textbf{09:00} budzik, wstajesz, odsłaniasz okno.
\end{itemize}
Rachunek: okno to $12 \times 60 = 720$ minut. Zasypianie zabiera 30, a nocne przebudzenie 30, więc zostaje 660 minut w łóżku, z czego sen to 480 minut (8 godzin). Pozostaje 180 minut, czyli trzy godziny, które nie są snem, mimo że są w łóżku. Sen zajmuje w oknie $480/720 \approx 67\%$.
\end{przyklad}

Ten wynik nie jest sam w sobie zły. Pokazuje jednak, że liczba 12 jest deklaracją o tym, ile czasu chcesz mieć w łóżku, a nie o tym, ile będziesz spać. \verify{Typowa wydajność snu u zdrowego dorosłego, czyli stosunek czasu przespanego do czasu spędzonego w łóżku, bywa podawana jako wysoka, powyżej osiemdziesięciu procent.} Jeśli Twoja wydajność jest wyraźnie niższa, nie wydłużaj okna, tylko przyjrzyj się, co zajmuje te trzy godziny.

\section{Światło: rano jasno, wieczorem ciemno}

Rytm dobowy zsynchronizowany jest głównie przez światło. Rano jego działanie jest dla rytmu korzystne, bo sygnalizuje początek doby. Praktycznie oznacza to wyjście na zewnątrz w ciągu pierwszej godziny po wstaniu, nawet na kilka minut, nawet w pochmurny dzień. Zamknięte pomieszczenie przy sztucznym świetle to dla organizmu inny sygnał niż ulica o dziewiątej rano. \verify{Ekspozycja na światło dzienne rano wspiera wcześniejsze ustawienie rytmu dobowego, a światło przez szybę daje wyraźnie słabszy efekt niż na zewnątrz.}

Wieczór działa odwrotnie. Jasne, białe światło w ostatnich godzinach przed snem opóźnia sygnał \glqq{}już noc\grqq{}. Zasada robocza: od 19:00 przygaszasz główne oświetlenie, a od 20:00 używasz lampek z ciepłym światłem. Przy porze snu o 21:00 to dwie godziny półmroku. Nie chodzi o ciemność absolutną, tylko o to, żeby wieczór nie wyglądał jak dzień.

\section{Kofeina i alkohol}

Kofeina działa długo. Jej czas półtrwania, czyli czas, po którym z organizmu znika połowa dawki, \verify{wynosi u wielu dorosłych około pięciu godzin, przy dużym rozrzucie między osobami}. Z tego prostego faktu wynika rachunek, który warto zrobić raz. Kawa o 14:00 i pora snu o 21:00 to siedem godzin. Przy półtrwaniu pięciu godzin zostaje $0{,}5^{7/5} \approx 0{,}38$, czyli około 38\% początkowej dawki. Kawa o 16:00 zostawia po pięciu godzinach połowę, a kawa o 10:00 po jedenastu godzinach zostawia około 22\%. Dwie kawy po obiedzie to więc nie dwie drobne decyzje, tylko dwie dawki, które wciąż pracują o północy.

Reguła robocza, a nie wynik badania: ostatnia kawa najpóźniej o 14:00, jeśli pora snu to 21:00. Przy zmęczeniu po obiedzie lepiej zrobić krótką przerwę albo spacer niż sięgnąć po kolejną dawkę. Jeśli nie wiesz, czy kofeina wpływa na Twój sen, odstaw ją na tydzień i zanotuj zasypianie.

Alkohol wygląda na sojusznika, bo wieczorem ułatwia zaśnięcie. Druga połowa nocy wygląda inaczej. \verify{Alkohol przed snem zwykle skraca pierwszą fazę zasypiania, a potem fragmentuje sen w jego drugiej części, kiedy organizm zaczyna go rozliczać.} Dlatego wieczorny kieliszek w tej książce jest traktowany jak wydatek z przyszłego budżetu, a nie jak lek.

\begin{falle}[Kieliszek na zaśnięcie]
Wino o 21:00 ma pomóc zasnąć i rzeczywiście często tak działa. Problem pojawia się o czwartej nad ranem: budzisz się, nie możesz wrócić do snu i uznajesz, że za mało śpisz, więc następnym razem bierzesz kieliszek ponownie. Korekta: mierz to, co się dzieje po przebudzeniu, a nie tylko godzinę zaśnięcia. Jeśli sen rozpada się w drugiej połowie okna, a wieczorem pił alkohol, to pierwsza podejrzana jest właśnie ta sprawa.
\end{falle}

\section{Sypialnia}

Sypialnia ma trzy zadania: ma być ciemna, cicha i chłodna. Temperatura jest najłatwiejsza do zmierzenia i najczęściej zaniedbywana. \verify{Temperatura w sypialni uznawana za wygodną do snu mieści się zwykle w przedziale od kilkunastu do dziewiętnastu stopni Celsjusza.} Termometr kosztuje mniej niż jeden wieczór w restauracji, a odczyt przez tydzień powie więcej niż ogólna zasada.

Ciemność oznacza zasłony, które zatrzymują uliczną latarnię, oraz brak diod świecących w ciemności. Cisza nie oznacza zupełnej ciszy, tylko brak nagłych dźwięków. Szum wentylatora albo białego szumu bywa bardziej pomocny niż cisza, bo zagłusza przypadkowe odgłosy. \verify{Wpływ hałasu nocnego na fragmentację snu jest udokumentowany, ale siła tego efektu zależy od osoby i natężenia dźwięku.}

Ostatnia zasada to przeznaczenie pomieszczenia. Łóżko do spania, nie do pracy, jedzenia i oglądania ekranów. Gdy łóżko ma tylko jedną funkcję, organizm łatwiej kojarzy je z zasypianiem. To reguła ergonomiczna, nie medyczna: nie wymaga żadnej liczby, tylko konsekwencji.

\section{Rutyna wieczorna i elektronika}

Rutyna wieczorna ma być krótka i powtarzalna. Czterdzieści pięć minut wystarcza, żeby oddzielić dzień od nocy, i mieści się w budżecie bez żadnego nadrabiania. W przykładzie z poprzedniego punktu zaczyna się o 20:15, a kończy o 21:00:

\begin{itemize}
\item \textbf{20:15--20:30} prysznic lub ciepła kąpiel (15 minut). \verify{Ciepła kąpiel na godzinę lub dwie przed snem bywa związana z łatwiejszym zasypianiem, a wyniki badań są niejednoznaczne.}
\item \textbf{20:30--20:40} plan na jutro na kartce: trzy rzeczy, ubranie, torba (10 minut).
\item \textbf{20:40--21:00} czytanie papierowej książki przy lampce (20 minut).
\end{itemize}

Suma wynosi $15 + 10 + 20 = 45$ minut. Plan na jutro jest tu najważniejszy, bo zamyka pętlę myśli, która w przeciwnym razie krąży po łóżku. Telefon ładujesz poza sypialnią od 20:15, nie w nocnym stoliku. Ekrany w ostatniej godzinie przed snem to temat, o którym mówi się dużo i często z przesadą. Konkretna zasada: żadnych wiadomości, których nie da się zostawić do rana. \verify{Wpływ ekranów na zasypianie zależy od jasności, treści i indywidualnej wrażliwości; dowody na wielkość tego efektu są mieszane.}

\section{Za dużo łóżka, drzemki i przesuwanie pory}

Dwanaście godzin to górna granica budżetu, a nie cel. Jeśli śpisz dwanaście godzin i nadal jesteś zmęczony, problem zwykle nie leży w tym, że za mało czasu, tylko w tym, że czas jest źle wykorzystany. Długie przebywanie w łóżku bez potrzeby snu zwykle skraca i rozdrabnia sen: pojawiają się płytkie fazy, wybudzenia, leżenie na granicy snu. \verify{Dłuższe okno w łóżku u osób, które nie potrzebują tyle snu, wiąże się z gorszą jakością snu, a nie z jego pogłębieniem.} Wtedy właściwym ruchem jest skrócenie okna o kwadrans, nie dokładanie kolejnego.

Kiedy 12 godzin to za dużo, sygnały są praktyczne: budzisz się przed budzikiem i nie możesz wrócić do snu, leżysz dłużej niż piętnaście minut bez zasypiania, w ciągu dnia nie czujesz się wyspany mimo długiego okna. To nie są diagnozy, tylko powody, żeby wrócić do rozdziału o potrzebach snu i ewentualnie porozmawiać z lekarzem.

Drzemki są osobnym tematem. Krótka drzemka po południu może poprawić samopoczucie. Długa lub późna zmienia rytm, bo zabiera sen z wieczoru. Reguła robocza: drzemka do około dwudziestu minut, najpóźniej o 15:00. \verify{Drzemki do dwudziestu minut są zwykle opisywane jako najmniej szkodliwe dla snu nocnego.}

\subsection{Przesuwanie pory o kwadranse}

Jeśli chcesz przesunąć porę snu o godzinę wcześniej, nie rób tego jednym skokiem. Skok o godzinę zwykle kończy się leżeniem z otwartymi oczami. Krok o 15 minut jest mały na tyle, że organizm go przyjmuje, a rachunek jest prosty: 60 minut dzieli się przez 15, czyli cztery kroki. Przy jednym kroku dziennie przesuwasz się o godzinę w cztery dni. Przy 23:00 jako punkcie wyjścia i 21:00 jako celu: $120 / 15 = 8$ kroków, czyli osiem dni. Nie zmieniasz pobudki w trakcie, dopóki nie wyczerpiesz kroków zasypiania. Pobudka jest kotwicą, a nie drugą zmienną do przesuwania.

Ta zasada wymaga cierpliwości, a nie siły woli. Jeśli jakiś krok nie działa przez dwa dni, zostajesz na nim, zamiast cofać się o godzinę.

\begin{uwaga}
Ta część książki opisuje higienę rytmu, a nie leczenie bezsenności. Jeśli problemy z zasypianiem lub z utrzymaniem snu trwają tygodniami, a dzień jest przez to trudny, to pytanie do lekarza, nie do kolejnego kroku o kwadrans.
\end{uwaga}

\begin{cwiczenia}
\exitem{Budzik masz ustawiony na 6:30, a okno snu ma mieć dwanaście godzin. O której godzinie kładziesz się do łóżka?}{Od 6:30 odejmujemy 12 godzin i dostajemy 18:30. To wczesna pora, więc warto sprawdzić, czy da się jej dotrzymać w praktyce. Jeśli nie, zmniejsz okno albo przesuń pobudkę.}
\exitem{Śpisz osiem godzin w oknie dwunastogodzinnym. Ile godzin spędzasz w łóżku bez snu i jaki jest udział snu w oknie?}{Cztery godziny poza snem: zasypianie, budzenie się i czekanie. Sen zajmuje $8/12 \approx 67\%$ okna.}
\exitem{Ostatnia kawa o 15:30, pora snu o 22:30, czas półtrwania przyjęty na pięć godzin. Ile godzin dzieli kawę od snu i jaka część dawki zostaje?}{Siedem godzin. Zostaje $0{,}5^{7/5} \approx 0{,}38$, czyli około 38\% dawki. Wynik zależy od przyjętego czasu półtrwania, który w tekście jest oznaczony do sprawdzenia.}
\exitem{W tygodniu śpisz 21:00--09:00, w sobotę 23:00--11:00. O ile przesuwa się środek snu i dlaczego to ma znaczenie?}{Środek snu w tygodniu wypada o 3:00, w sobotę o 5:00, czyli przesunięcie wynosi dwie godziny. Znaczenie ma to, że rytm ustawia się według pobudki, a skok na weekend przesuwa go w inną stronę.}
\exitem{Chcesz przesunąć porę snu z 23:00 na 21:00, krokami po 15 minut, zmieniając jeden krok dziennie. Ile dni potrzebujesz?}{Różnica to 120 minut, czyli $120 / 15 = 8$ kroków. Przy jednym kroku dziennie to osiem dni, z zastrzeżeniem, że krok, który nie działa przez dwa dni, się powtarza.}
\exitem{Rutyna wieczorna ma trwać od 20:15 do 21:00. Kąpiel zajmuje 15 minut, plan na jutro 10, czytanie ma zająć resztę. Ile minut zostaje na czytanie i czy to się sumuje?}{Od 20:15 do 21:00 jest 45 minut. Kąpiel i plan zajmują 25, więc na czytanie zostaje 20 minut. Suma wynosi $15 + 10 + 20 = 45$.}
\exitem{Temperatura w sypialni wieczorem wynosi wyraźnie więcej, niż wygodnie się śpi. Podaj dwie zmiany, które kosztują najmniej.}{Przewietrzenie pokoju tuż przed pójściem spać i ograniczenie nagrzewania pomieszczenia w ciągu dnia, na przykład przez zasłonięcie okien od południa. Obie zmiany są tanie i odwracalne.}
\exitem{Budzisz się o 4:00, nie możesz zasnąć, a w ciągu dnia jesteś zmęczony mimo dwunastu godzin w łóżku. Które dwa elementy tego rozdziału sprawdzisz najpierw?}{Długość okna, czyli czy leżysz za długo bez snu, oraz wieczorny alkohol lub kofeinę po południu. Oba są szybkie do sprawdzenia przez tydzień zapisu.}
\end{cwiczenia}

\zadanie{Przez siedem dni wstawaj o tej samej godzinie, z tolerancją 15 minut, także w weekend. Każdego ranka zapisz godzinę pobudki, godzinę zaśnięcia (z przybliżeniem) i to, czy wyszedłeś na zewnątrz w ciągu godziny. Miarą sukcesu jest to, że w siedmiu dniach pobudka mieści się w przedziale 15 minut, a w tabeli nie ma ani jednej luki.}
